\subsection*{Intended Learning Outcomes - Lecture 09}

The student should be able to:
\begin{itemize}
\item know the definition of marginalization of SCMs;
\item know the basic commutation and compatibility properties of basic operations on SCMs (marginalization, twinning, intervening, preservation of unique solvability / simplicity / equivalence);
\item be able to construct a marginalized SCM;
\item know the definition of the graph of an SCM;
\item be able to draw the graph of a given SCM, or to construct an SCM with a given graph;
\item be able to twin a graph;
\item know the basic compatibility relations between the operation that maps an SCM to a graph and other basic SCM operations;
\item know the definition of acyclic SCM;
\item understand and appreciate the differences between different causal model classes;
\item know the definitions of acyclification of a graph and of an SCM, and their relation;
\item know the relation between $\sigma$-separation and $d$-separation on an acyclification;
\item know the global Markov property for simple SCMs;
\item be able to apply the global Markov property for simple SCMs;
\item write down the 3 do-calculus rules for simple SCMs and reason with them;
\item apply the 3 do-calculus rules for simple SCMs in an iterative manner on (small, not too complicated) graphs for the identification of causal effects (i.e.\ reducing the number $\doit$-terms);
\item state and apply the back-door covariate adjustment criterion and formula for simple SCMs;
\end{itemize}
